\section{PHM2012 结果}

\subsection{Official Split 与分析 Caveat}

PHM2012 的主结果来自 Step J，即 \config{phm2012_official} + \config{manual_basic}。

Step K 在同一 official split 上运行 \config{manual_tsfresh_basic}，用于对照 tsfresh minimal 是否增加新信息。

本节只整理已有归档报告，不增加新实验。

PHM2012 的 train/val/test 按轴承隔离。

训练集包含 Bearing1/2/3 工况下的官方训练轴承。

验证集使用 Bearing1\_3 与 Bearing2\_3。

测试集使用剩余官方测试轴承。

Health State 与 Early Fault 的解释需要谨慎。

它们依然是 HI/FPT 派生 pseudo-label，不是人工逐点标注的真实状态。

因此，本节对这两个任务强调“分布可分性”和“FPT 前后变化”，不把它们等同为真实维修事件。

\begin{table}[htbp]
\centering
\caption{PHM2012 manual\_basic 的任务级 top evidence}
\label{tab:phm2012_feature_evidence}
\scriptsize
\begin{tabularx}{\textwidth}{llrY}
\toprule
任务 & 特征 & score/rank & 解释 \\
\midrule
RUL & \feature{h__time__mean_abs} & rank 1, score 0.544 & Pearson=-0.638, Spearman=-0.659；水平幅值随退化增长。\\
RUL & \feature{mag__time__mean} & rank 2, score 0.537 & magnitude amplitude 支持 late-life growth。\\
RUL & \feature{h__time__rms} & rank 4, score 0.527 & 与 \feature{h__time__std} 冗余，可作为 secondary。\\
Health State & \feature{h__time__mean_abs} & rank 1, score 0.944 & MI=0.486, Fisher=0.789, ANOVA F=1979.4；boxplot 分离强。\\
Health State & \feature{h__time__std} & rank 2, score 0.854 & horizontal dispersion 随伪健康阶段上升。\\
Early Fault & \feature{h__time__mean_abs} & rank 1, score 0.994 & AUC=0.836, Cohen's d=1.000；FPT 后幅值上升。\\
Early Fault & \feature{mag__time__mean} & rank 2, score 0.983 & AUC=0.824, Cohen's d=1.012；magnitude pre/post FPT 可分。\\
Reference & \feature{mag__time__rms} & RUL rank 6 & 实际 HI/FPT source，仅作参考。\\
\bottomrule
\end{tabularx}
\end{table}


\subsection{RUL 特征分析}

PHM2012 RUL 的 top feature 是 \feature{h__time__mean_abs}。

它的 RUL score 为 0.544。

Pearson 为 -0.638，Spearman 为 -0.659，Kendall 为 -0.516。

这说明水平通道平均绝对幅值与归一化 RUL 有稳定负相关。

负相关的含义是：随着样本接近失效，RUL 下降，而振动幅值上升。

\feature{mag__time__mean} 和 \feature{mag__time__mean_abs} 并列 rank 2。

它们的 RUL score 为 0.537。

这两个 magnitude 特征和 XJTU-SY 的主结论一致，说明幅值/能量 family 在跨数据集上是共同模式。

\feature{h__time__rms} 和 \feature{h__time__std} 分别 rank 4 和 rank 5。

它们与 \feature{h__time__mean_abs} 描述同一 horizontal amplitude/dispersion story。

因此它们适合放入 B 级或 compact RUL subset 的 secondary group。

频域特征不是 PHM2012 RUL 主线。

\feature{v__spectral__rms_frequency} 和 \feature{v__spectral__peak_frequency} 的 RUL 相关性明显弱于幅值类特征。

这说明当前 PHM2012 RUL 主要由振动幅值随退化增加表达，而不是由频谱中心或峰值频率主导。

\begin{figure}[H]
\centering
\includegraphics[width=0.92\textwidth]{phm2012/manual/rul_top_features.png}
\caption{PHM2012 official split 下 RUL 任务的 top features。}
\label{fig:phm2012_rul_top_features}
\end{figure}

图 \ref{fig:phm2012_rul_top_features} 显示 horizontal 和 magnitude amplitude 特征占据前列。

读图时应先看非 label-source 特征。

\feature{h__time__mean_abs} 是最重要的独立候选。

\feature{mag__time__mean} 和 \feature{mag__time__mean_abs} 提供 magnitude 方向的支持。

\feature{mag__time__rms} 若出现在前列，只能作为 Reference。

这张图支撑 PHM2012 RUL 的结论：以 horizontal amplitude 为第一候选，magnitude amplitude 为共同退化信号。

\subsection{Health State 特征分析}

PHM2012 Health State strongly favors horizontal-channel amplitude features。

\feature{h__time__mean_abs} 的 health score 为 0.944，rank 1。

它的 mutual information 为 0.486，Fisher score 为 0.789，ANOVA F 为 1979.4。

这三个指标共同说明，\feature{h__time__mean_abs} 对伪健康阶段有明显分离能力。

\feature{h__time__std} 和 \feature{h__time__rms} 分别 rank 2 和 rank 3。

它们的 boxplot 形态接近，说明 horizontal dispersion 和 energy 是同一退化现象的不同统计表达。

\feature{mag__time__mean} 和 \feature{mag__time__mean_abs} 是 B 级支持。

它们不如 horizontal top features 强，但仍有稳定分离能力。

\feature{mag__spectral__entropy} 进入 C 级。

它说明频谱复杂度在 PHM2012 Health State 中有辅助信号，但不是主线。

\begin{figure}[H]
\centering
\includegraphics[width=0.92\textwidth]{phm2012/manual/health_state_boxplots.png}
\caption{PHM2012 official split 下 Health State 的箱线图。}
\label{fig:phm2012_health_boxplots}
\end{figure}

图 \ref{fig:phm2012_health_boxplots} 的重点是水平通道幅值类特征的状态分布。

如果健康阶段越严重，中位数和箱体整体上移，就说明该特征与退化阶段一致。

PHM2012 中，\feature{h__time__mean_abs}、\feature{h__time__std}、\feature{h__time__rms} 具有清晰的状态分离。

这和 XJTU-SY Step F 的 horizontal pattern 一致。

但 PHM2012 没有像 XJTU-SY 那样做跨工况 train=one-condition/test=another-condition 的压力测试。

所以最终写法应是：PHM2012 official split 支持 horizontal amplitude Health State；跨数据集共同主线仍是 amplitude family。

\subsection{Early Fault 特征分析}

PHM2012 Early Fault 同样由 amplitude family 主导。

\feature{h__time__mean_abs} rank 1，early fault score 为 0.994。

它的 AUC 为 0.836，Cohen's d 为 1.000，mean shift 为 0.220。

这说明 FPT 后水平平均绝对幅值明显上升。

\feature{mag__time__mean} 和 \feature{mag__time__mean_abs} 并列 rank 2。

它们的 AUC 为 0.824，Cohen's d 为 1.012，mean shift 为 0.317。

这说明 magnitude amplitude 对 FPT 前后分离同样有效。

\feature{h__time__std} 和 \feature{h__time__rms} 是 strong secondary。

它们捕捉的是水平通道离散度和能量增加。

这和早期故障中振动冲击/幅值扩大的物理直觉一致。

频域特征如 \feature{v__spectral__rms_frequency} 和 \feature{v__spectral__peak_frequency} 出现在 C 级。

它们有一定 pre/post FPT 分离能力，但不如 amplitude family 稳定。

\begin{figure}[H]
\centering
\includegraphics[width=0.92\textwidth]{phm2012/manual/early_fault_effects.png}
\caption{PHM2012 official split 下 Early Fault 的特征效应图。}
\label{fig:phm2012_early_fault_effects}
\end{figure}

图 \ref{fig:phm2012_early_fault_effects} 需要看 FPT 前后均值差和分布重叠。

PHM2012 中，horizontal/magnitude amplitude 特征在 FPT 后有明显上升。

这支撑把 \feature{h__time__mean_abs}、\feature{mag__time__mean}、\feature{mag__time__mean_abs} 放入 A 级。

同时，图也提醒我们 Early Fault 是 pseudo-label。

如果某个 Reference 特征强，只能说明它和 HI/FPT 规则一致。

非 label-source 的 horizontal/magnitude amplitude 特征才是独立候选。

\subsection{\config{manual_tsfresh_basic} 对照}

Step K 在 PHM2012 上成功运行 full-size \config{manual_tsfresh_basic}。

该 ranking 中有 44 个 ranked features，其中 20 个为 \feature{tsfresh__} 特征。

两个长度特征 \feature{tsfresh__h__length} 和 \feature{tsfresh__v__length} 因为常量被清理掉。

\begingroup
\scriptsize
\setlength{\tabcolsep}{3pt}
\renewcommand{\arraystretch}{1.18}
\begin{longtable}{@{}p{2.1cm}p{1.3cm}p{5.2cm}p{5.4cm}@{}}
\caption{PHM2012 manual\_tsfresh\_basic 对照}
\label{tab:phm2012_tsfresh_comparison}\\
\toprule
任务 & 数量 & 代表 tsfresh 特征 & 解释与决策 \\
\midrule
\endfirsthead
\toprule
任务 & 数量 & 代表 tsfresh 特征 & 解释与决策 \\
\midrule
\endhead
RUL & 5 & \feature{tsfresh__h__root_mean_square}; \feature{tsfresh__h__standard_deviation}; \feature{tsfresh__h__variance} & 重复 RMS/std/variance amplitude 信息；不改变主线。\\
Health State & 4 & \feature{tsfresh__h__standard_deviation}; \feature{tsfresh__h__root_mean_square}; \feature{tsfresh__h__maximum} & 确认 horizontal amplitude story；没有新增物理特征族。\\
Early Fault & 4 & \feature{tsfresh__h__standard_deviation}; \feature{tsfresh__v__standard_deviation}; \feature{tsfresh__v__root_mean_square} & 支持幅值/离散度 pre/post FPT 变化；保留为后续工程候选，不进入当前主线。\\
\bottomrule
\end{longtable}
\endgroup


tsfresh 对照的关键不是“是否进入 top10”，而是“进入 top10 的东西是否提供新信息”。

RUL 中进入 top10 的 tsfresh 特征包括：

\begin{itemize}[leftmargin=2em]
  \item \feature{tsfresh__h__root_mean_square}
  \item \feature{tsfresh__h__standard_deviation}
  \item \feature{tsfresh__h__variance}
\end{itemize}

这些分别对应手工 RMS、standard deviation 和 variance。

Health State 中 \feature{tsfresh__h__maximum} 和 \feature{tsfresh__h__absolute_maximum} 提供的是 horizontal amplitude extreme。

Early Fault 中 vertical tsfresh standard deviation/root mean square 与手工 vertical dispersion/RMS 重复。

因此 Step K 的结论是：PHM2012 tsfresh 成功运行，但其有效信息主要重复 amplitude/statistical family。

\begin{figure}[H]
\centering
\includegraphics[width=0.92\textwidth]{phm2012/tsfresh/feature_recommendation_matrix.png}
\caption{PHM2012 \config{manual_tsfresh_basic} 的特征推荐矩阵。}
\label{fig:phm2012_tsfresh_recommendation_matrix}
\end{figure}

图 \ref{fig:phm2012_tsfresh_recommendation_matrix} 显示部分 \feature{tsfresh__} 特征进入 RUL、Health State 和 Early Fault 的推荐区域。

但是这些特征集中在 RMS、standard deviation、variance、maximum。

它们确认了 PHM2012 的 amplitude-dominant story。

它们没有提供足以替换 \config{manual_basic} 的新物理特征族。

\subsection{PHM2012 最终决策}

\begin{table}[htbp]
\centering
\caption{PHM2012 三任务推荐特征摘要}
\label{tab:phm2012_top_features}
\small
\begin{tabularx}{\textwidth}{lYYY}
\toprule
任务 & A 级 & B 级 & C 级或 Reference \\
\midrule
RUL & \feature{h__time__mean_abs}, \feature{mag__time__mean}, \feature{mag__time__mean_abs} & \feature{h__time__rms}, \feature{h__time__std}, \feature{v__time__mean_abs}, \feature{mag__time__std} & Reference: \feature{mag__time__rms} \\
Health State & \feature{h__time__mean_abs}, \feature{h__time__std}, \feature{h__time__rms} & \feature{mag__time__mean}, \feature{mag__time__mean_abs} & C: \feature{mag__spectral__entropy}; Reference: \feature{mag__time__rms} \\
Early Fault & \feature{h__time__mean_abs}, \feature{mag__time__mean}, \feature{mag__time__mean_abs} & \feature{h__time__std}, \feature{h__time__rms}, \feature{v__time__mean_abs}, \feature{v__time__std} & C: \feature{v__spectral__rms_frequency}, \feature{v__spectral__peak_frequency}; Reference: \feature{mag__time__rms} \\
\bottomrule
\end{tabularx}
\end{table}


PHM2012 的主线特征集仍为 \config{manual_basic}。

RUL 使用 \feature{h__time__mean_abs} 与 magnitude amplitude。

Health State 使用 horizontal amplitude family。

Early Fault 使用 horizontal/magnitude amplitude，并把 spectral frequency 作为 C 级 secondary。

\config{manual_tsfresh_basic} 可运行，但当前不作为主线。

后续如果要把 tsfresh 纳入 baseline，需要先证明它带来的不是 RMS/std/variance 的重复，而是对模型泛化或解释性有实际增益。
